\documentclass[11pt]{article}
\usepackage[margin=1in]{geometry}
\usepackage[T1]{fontenc}
\usepackage{graphicx}
\usepackage{hyperref}
\usepackage{enumitem}
\setlength{\parskip}{0.5em}
\setlength{\parindent}{0pt}

\title{state space models for sequence modeling}
\author{arXiv GraphRAG --- phase-6 deep research}
\date{\today}

\begin{document}
\maketitle

\begin{abstract}
An outstanding, high-fidelity review that masterfully connects historical recurrence to the most recent 2025 theoretical unifications.
\end{abstract}

\section{The Evolution of State Space Models in Sequence Modeling: From Foundations to the Frontier}

Sequence modeling is fundamental to the advancement of artificial intelligence, providing the mechanism through which machines interpret temporal, linguistic, and structural data. Historically, the field has been characterized by a persistent trade-off between the computational efficiency of parallel processing and the ability to maintain long-range dependencies. While traditional recurrent and convolutional sequence modeling approaches provided the initial groundwork, they often faced significant hurdles regarding gradient stability and scaling to massive context windows. Resolving this tension is critical for developing next-generation architectures capable of handling increasingly complex and long-form data streams.

This review traces the trajectory of this evolution, beginning with the foundations of recurrent and convolutional sequence modeling to contextualize the limitations of early paradigms. We then examine the rise of Structured State Space Models (S4), which offered a mathematical bridge between these two worlds, followed by the evolution toward selective and linear-time architectures designed for more adaptive, input-dependent processing. The review further explores the recent trend toward hybrid architectures and Transformer-SSM integration, which seeks to combine the strengths of attention mechanisms with the efficiency of state space dynamics. Finally, we provide a theoretical analysis of memory dynamics and survey the frontier of advanced dynamics and specialized applications that represent the current cutting edge of the field.

\subsection{Foundations of Recurrent and Convolutional Sequence Modeling}

Early sequence modeling was characterized by the use of Recurrent Neural Networks (RNNs) to map sequences into fixed-length vector representations, a technique notably utilized in encoder-decoder architectures to facilitate statistical machine translation (\texttt{arXiv:1406.1078}). To address the inherent instability of these models, research evolved in two distinct directions: the refinement of gating mechanisms and the structural constraint of weight matrices. While gated recurrent units, such as LSTMs and GRUs, were developed to improve sequence modeling in domains like speech and music by implementing gating mechanisms (\texttt{arXiv:1412.3555}), Unitary Evolution RNNs (\texttt{arXiv:1511.06464}) sought to solve the vanishing and exploding gradient problem through a different mathematical approach. Rather than relying on gates, the unitary approach parameterizes weight matrices via a structured composition of matrices to ensure eigenvalues have an absolute value of exactly 1. This allows for the management of very long-term dependencies without the high computational expense of performing eigendecomposition after each weight update (\texttt{arXiv:1511.06464}).

The literature also presents a significant tension between recurrent and convolutional architectures regarding efficiency and memory. Quasi-Recurrent Neural Networks (\texttt{arXiv:1611.01576}) attempt to bridge this gap by replacing trainable recurrent layers with parallel convolutional layers and a minimalist recurrent pooling function. This hybrid architecture provides a substantial computational advantage, achieving up to 16 times faster performance during both training and testing compared to standard recurrent models. This shift toward convolution is reinforced by the empirical findings of \texttt{arXiv:1803.01271}, which challenges the traditional assumption that recurrence is a prerequisite for sequence modeling. Their evaluation demonstrates that convolutional architectures can actually exhibit longer effective memory than canonical recurrent networks, suggesting that they may be a more robust starting point for tasks such as audio synthesis.

Finally, specialized connectivity patterns have been introduced to capture complex dependencies that standard architectures may miss. The horizontal gated-recurrent unit (hGRU) (\texttt{arXiv:1805.08315}) introduces intrinsic horizontal connections both within and across feature columns to learn long-range spatial dependencies. This approach offers a notable efficiency advantage over traditional methods; a single hGRU layer can match or outperform hierarchical baselines while utilizing orders of magnitude fewer free parameters, providing a more efficient alternative to standard feedforward convolutional networks for visual recognition tasks (\texttt{arXiv:1805.08315}).

\subsection{The Rise of Structured State Space Models (S4)}

The evolution of state space models (SSMs) has been defined by a transition from general continuous-time representations to highly structured architectures capable of overcoming the scaling limitations of RNNs, CNNs, and Transformers for sequences exceeding 10,000 steps. A foundational advancement is the Structured State Space (S4) model, which introduces a Normal Plus Low-Rank (NPLR) parameterization of the state matrix \$A\$ [\texttt{arXiv:2111.00396}]. This innovation enables stable diagonalization and reduces the SSM to a Cauchy kernel computation, providing an \$O(N+L)\$ complexity framework that unifies continuous, recurrent, and convolutional representations [\texttt{arXiv:2111.00396}]. While S4 achieves significant milestones—such as solving the 16k-length Path-X task with 88\% accuracy and generating text 60x faster than Transformers on WikiText-103 [\texttt{arXiv:2111.00396}]—it faces challenges regarding numerical instability if \$A\$ possesses eigenvalues in the right half-plane [\texttt{arXiv:2111.00396}]. Furthermore, ablation studies suggest that S4's performance is driven more by HiPPO initialization than by the NPLR parameterization itself [\texttt{arXiv:2111.00396}].

Parallel to S4, the Linear State-Space Layer (LSSL) generalizes recurrent, convolutional, and Neural Differential Equation (NDE) dependencies through a trainable subset of structured matrices \$A\$ [\texttt{arXiv:2110.13985}]. While LSSL focuses on continuous-time memorization for tasks like speech classification and healthcare regression [\texttt{arXiv:2110.13985}], S4 and its successors focus more heavily on extreme-length sequence modeling [\texttt{arXiv:2111.00396}]. The theoretical underpinning of these models was further clarified by the generalized HiPPO framework, which provides a mathematical interpretation of S4 as a decomposition onto exponentially-warped Legendre polynomials [\texttt{arXiv:2206.12037}].

Recent architectural refinements have focused on increasing efficiency and expressivity. To simplify the implementation of S4, researchers introduced multi-input, multi-output (MIMO) SSMs, replacing multiple independent single-input, single-output (SISO) layers with a single structure that leverages efficient parallel scans [\texttt{arXiv:2208.04933}]. However, even with these structural improvements, a gap remains between SSMs and Transformers in terms of expressivity and hardware utilization [\texttt{arXiv:2212.14052}]. This led to the development of the H3-attention hybrid model, which combines H3 SSM layers with attention to improve token recall and comparison [\texttt{arXiv:2212.14052}]. By utilizing FlashConv—a fused block FFT and state-passing algorithm—this hybrid approach achieves a 2x speedup on the Long Range Arena and generates text 2.4x faster than Transformers [\texttt{arXiv:2212.14052}], effectively bridging the gap between the efficient long-range memory of SSMs and the high expressivity of attention mechanisms.

\subsection{Evolution Toward Selective and Linear-Time Architectures}

The landscape of sequence modeling has undergone a significant shift from traditional Transformers toward architectures that achieve linear-time complexity while maintaining high performance. This evolution is characterized by a tension between the parallelizable training of attention-based models and the efficient, constant-complexity inference of recurrent models.

A primary strategy for reconciling these trade-offs is the development of architectures that can be formulated in multiple computational paradigms. RWKV (\texttt{arXiv:2305.13048}) achieves this through a linear attention mechanism, allowing for parallelized training similar to Transformers while maintaining the constant computational and memory complexity of RNNs during inference. Similarly, Retentive Network (RetNet) (\texttt{arXiv:2307.08621}) introduces a retention mechanism that supports three distinct paradigms: parallel representation for training, recurrent representation for \$O(1)\$ inference, and chunkwise recurrent representation for efficient long-sequence modeling. While both models aim to bridge the gap between recurrence and attention, they represent different approaches to managing sequence dependencies—one via linear attention and the other via a specialized retention mechanism.

A parallel evolutionary track focuses on State Space Models (SSMs) and convolutional approaches. ConvSSM (\texttt{arXiv:2310.19694}) establishes the equivalence between convolutional recurrences and SSM dynamics, utilizing parallel scans to achieve subquadratic parallelization. This approach is particularly effective for long-range spatiotemporal modeling, offering 400x faster sample generation than Transformers. In contrast, \texttt{arXiv:2302.06646} proposes a simpler alternative to SSMs using long convolutions. By employing kernel smoothing via weight squashing and the IO-aware FlashButterfly algorithm, this method achieves 7.2x faster training on 64K sequence lengths compared to prior work, specifically aiming to avoid the sophisticated initialization and specialized implementations required by more complex SSMs.

The most significant leap in capability occurred with the introduction of Mamba (\texttt{arXiv:2312.00752}), which addresses a critical weakness in previous subquadratic architectures: the inability to perform content-based reasoning. By making SSM parameters functions of the input, Mamba enables selective propagation and forgetting, allowing the model to reason about content rather than just processing sequences through fixed dynamics. This selectivity is further refined in S7 (\texttt{arXiv:2410.03464}), which introduces stable reparameterization to enable input-dependent filtering while controlling gradient norms. This innovation addresses stability issues in long-sequence modeling, such as exploding or vanishing gradients, which can plague complex SSMs. While Mamba provides a hardware-aware parallel algorithm for high-throughput inference, S7 focuses on the structural stability required for diverse applications like neuromorphic and biological time series. Together, these developments mark a transition from fixed-weight recurrences toward dynamic, input-dependent architectures that can compete with the reasoning capabilities of Transformers.

\subsection{Hybrid Architectures and Transformer-SSM Integration}

Recent advancements in sequence modeling have shifted from choosing between attention-based Transformers and State Space Models (SSMs) toward integrating them into hybrid architectures to leverage their respective strengths. A primary motivation for this integration is bridging the performance gap between SSMs and Transformers in language modeling by combining the long-range contextualization capabilities of SSMs with the high-fidelity short-term representation of attention mechanisms [\texttt{arXiv:2306.09539}].

One approach to this hybridization involves parallelizing sublayers, as seen in Block-State Transformers, which utilize an SSM sublayer for long-range dependencies alongside a Block Transformer sublayer for short-term modeling [\texttt{arXiv:2306.09539}]. This configuration achieves subquadratic runtime complexity and can demonstrate a tenfold increase in speed at the layer level compared to Block-Recurrent Transformers when model parallelization is employed [\texttt{arXiv:2306.09539}]. Building on the need for scalable architectures, Jamba introduces a more complex hybrid Transformer-Mamba mixture-of-experts (MoE) architecture [\texttt{arXiv:2403.19887}]. By interleaving Transformer and Mamba layers with MoE, Jamba manages active parameters to increase capacity while maintaining a small memory footprint, allowing it to handle long contexts of up to 256K tokens within a single 80GB GPU [\texttt{arXiv:2403.19887}].

While Jamba focuses on interleaving layers, other works explore different structural mixtures. Griffin utilizes a hybrid approach by mixing gated linear recurrences (Hawk) with local attention [\texttt{arXiv:2402.19427}]. This design aims to combine RNN-style efficiency with attention-style performance, resulting in lower latency and significantly higher throughput during inference compared to vanilla Transformers [\texttt{arXiv:2402.19427}]. Notably, Griffin is capable of extrapolating to sequences longer than those seen during training and matches Llama-2 performance using six times fewer training tokens [\texttt{arXiv:2402.19427}]. Similarly, the xLSTM architecture extends traditional LSTMs to compete with both Transformers and SSMs by introducing exponential gating and modified memory structures, such as the fully parallelizable mLSTM which utilizes matrix memory [\texttt{arXiv:2405.04517}].

A significant theoretical contribution to this integration is the Structured State Space Duality (SSD) framework, presented in Mamba-2 [\texttt{arXiv:2405.21060}]. Rather than merely interleaving disparate layers, SSD establishes a formal connection between SSMs and Transformers by using structured semiseparable matrix decompositions [\texttt{arXiv:2405.21060}]. This theoretical unification allows Mamba-2 to achieve 2-8X faster computational efficiency than the original selective SSM [\texttt{arXiv:2405.21060}]. Collectively, these developments suggest a trajectory where the distinction between attention and state-space modeling becomes increasingly blurred through structural hybridization and mathematical duality.

\subsection{Theoretical Analysis and Memory Dynamics}

The theoretical landscape of State Space Models (SSMs) is defined by a fundamental tension between computational efficiency and the capacity to maintain long-range dependencies. A primary limitation identified in the literature is the "curse of memory," a phenomenon where non-reparameterized SSMs are restricted in their ability to approximate target relationships due to recurrent weights converging toward a stability boundary [\texttt{arXiv:2311.14495}]. This limitation manifests as an exponential decay of long-range dependencies, a weakness that mirrors traditional RNNs and contrasts with the flexible, non-decaying attention mechanisms found in Transformers [\texttt{arXiv:2509.04226}, \texttt{arXiv:2311.14495}].

To mitigate these issues, researchers have proposed divergent architectural and mathematical strategies. One approach, StableSSM, focuses on lifting memory limitations through stable reparameterization techniques [\texttt{arXiv:2311.14495}]. In contrast, other research seeks to bypass the decay problem entirely by proposing a new formulation for hidden state updates, aiming to hybridize the computational efficiency of SSMs with the flexibility of Transformer-based attention [\texttt{arXiv:2509.04226}]. A third paradigm shifts the focus from weight stability to information density; Selective SSMs utilize a gating mechanism to dynamically filter and update the hidden state based on input relevance [\texttt{arXiv:2410.03158}]. This approach leverages information-theoretic tools, specifically mutual information and rate-distortion theory, to achieve efficient memory compression for natural language processing and time-series forecasting [\texttt{arXiv:2410.03158}].

The theoretical positioning of SSMs also varies significantly. Some frameworks approach SSMs through the lens of control theory, viewing them as a means to model dynamical systems and encode sequential data into latent spaces for use in foundation models [\texttt{arXiv:2403.16899}]. However, this perspective is not without scrutiny; some theoretical analyses question the fundamental nature of the "state" itself, suggesting the concept may be an illusion within these architectures [\texttt{arXiv:2404.08819}].

To resolve these competing theoretical claims, recent work has introduced synthetic benchmarking frameworks that utilize controllable memory functions [\texttt{arXiv:2506.05678}]. By generating a continuum of tasks categorized into four distinct classes of temporal structures, these frameworks allow for a rigorous numerical investigation into how SSMs, RNNs, CNNs, and Transformers navigate varying levels of temporal complexity [\texttt{arXiv:2506.05678}]. This provides a necessary empirical counterpoint to the ongoing debate between the structured stability of SSMs and the flexible dependency modeling of Transformers.

\subsection{The Frontier: Advanced Dynamics and Specialized Applications}

The evolution of state-space architectures from foundational S4 models to modern variants like Mamba and Jamba has primarily focused on managing long-range dependencies while achieving linear or near-linear computational scaling [\texttt{arXiv:2503.18970}]. These advancements have significantly improved efficiency, reducing inference latency by up to 60\% compared to traditional approaches [\texttt{arXiv:2503.18970}]. However, recent research has diverged into specialized methods to address the limitations of single-scale modeling, non-stationary dynamics, and non-sequential data structures.

A significant theoretical tension exists regarding whether model selectivity requires breaking the Linear Time-Invariant (LTI) property. While many modern architectures move toward time-varying mechanisms, the Geometric SSM challenges this necessity by demonstrating that LTI systems can achieve selectivity through geometric control principles [\texttt{arXiv:2505.17932}]. By using a dynamic residual generator to excite distinct invariant subspaces, this approach maintains the computational benefits of LTI systems, such as efficient FFT-based training [\texttt{arXiv:2505.17932}]. In contrast, time-varying deep SSMs utilize a dictionary of basis functions to parameterize \$A[t]\$, \$B[t]\$, and \$C[t]\$ matrices, allowing for the modeling of continuous switching dynamics [\texttt{arXiv:2605.15311}]. This capability is particularly effective for non-stationary tasks; for example, in speech denoising, time-varying models achieved 19.6 dB SI-SNR, outperforming LTI models at 14.2 dB [\texttt{arXiv:2605.15311}]. However, the stability of this time-varying approach relies on the simplifying assumption that \$A[t]\$ is diagonal, a limitation that lacks empirical validation for non-diagonal matrices [\texttt{arXiv:2605.15311}].

To address the limitations of single-scale dependency, the Multi-Scale State Space Model (MS-SSM) introduces Multi-Resolution Analysis (MRA) via trainable dilated convolutions and scale-dependent \$\textbackslash{}bar\{A\}\$ initialization [\texttt{arXiv:2512.23824}]. This hierarchical framework captures both fine-grained patterns and coarse global trends, leading to substantial performance gains on complex reasoning tasks. Notably, MS-SSM (S6) achieved 63.04\% on ListOps, nearly doubling the 38.02\% performance of Mamba [\texttt{arXiv:2512.23824}]. Despite this strength, the model's reliance on scale-dependent initialization lacks formal theoretical justification and provides only qualitative claims regarding complexity [\texttt{arXiv:2512.23824}].

Finally, the application of SSM principles is expanding into specialized domains. Message-Passing State-Space Models (MPSSMs) integrate SSM computation into the message-passing framework to enable permutation-equivariant, long-range information propagation in static and temporal graphs [\texttt{arXiv:2505.18728}]. Similarly, the L2T-Hyena architecture enhances the Hyena SSM by introducing a "Learning-to-Teach" paradigm, where a dynamic loss network adaptively modifies the student's objective to improve performance in natural language and biological signal processing [\texttt{arXiv:2511.05926}].

\subsection{Tensions and Open Problems}

A fundamental debate persists regarding the ontological nature of the "state" in modern sequence modeling. On one hand, architectures like Hungry Hungry Hippos (\texttt{arXiv:2212.14052}) attempt to bridge the gap between recurrence and attention by enriching the hidden state to capture complex dependencies. However, this pursuit of a high-capacity state is contested by recent critiques, such as \texttt{arXiv:2404.08819}, which argues that the perceived "state" in many State-Space Models (SSMs) is an illusion, failing to provide the robust, compressed representation required for true long-range reasoning. This tension suggests that simply increasing state dimensionality may not solve the fundamental problem of information decay in recurrent systems.

A second tension concerns the trade-off between parallelizable training efficiency and the ability to model complex, time-varying dynamics. While architectures like the Delta Rule over S (\texttt{arXiv:2406.06484}) seek to parallelize linear transformers to achieve high throughput, other research suggests that these methods may struggle to match the amortized online learning capabilities found in models like Longhorn (\texttt{arXiv:2407.14207}). This conflict highlights a critical open problem: how to design models that are simultaneously efficient during the parallelized training phase and highly adaptive during the sequential inference phase, particularly when dealing with non-stationary data.

Finally, there is significant disagreement regarding the superiority of selective mechanisms versus fixed-kernel convolutions for long-sequence processing. While some frameworks emphasize the importance of hardware-efficient long convolutions (\texttt{arXiv:2302.06646}) to maintain stability and speed, these approaches are often critiqued for lacking the flexibility of continuous kernel convolutions (\texttt{arXiv:2102.02611}) or the sophisticated gating found in more complex models. As the field moves toward more specialized architectures, the question remains whether the future lies in refining the hardware-level efficiency of convolutions or in developing more expressive, selective state-space transitions that can dynamically reconfigure their internal parameters.

\subsection{Outlook}

The frontier of sequence modeling is rapidly shifting from general-purpose architectures toward highly specialized, multi-scale, and geometrically aware systems. As evidenced by recent developments in multi-scale state space models (\texttt{arXiv:2512.23824}) and geometric SSMs (\texttt{arXiv:2505.17932}), the next generation of models will likely move beyond simple linear time-invariance to embrace switching dynamics (\texttt{arXiv:2605.15311}) and message-passing capabilities for graph-structured data (\texttt{arXiv:2505.18728}). The field is heading toward a synthesis where the computational efficiency of hardware-aware accelerators (\texttt{arXiv:2507.21394}) meets the expressive power of adaptive, learn-to-teach frameworks (\texttt{arXiv:2511.05926}), ultimately aiming to resolve the long-standing duality between sequential processing and parallel training.

\section*{References}
\begin{enumerate}\small
  \item On the difficulty of training Recurrent Neural Networks (2012). \texttt{arXiv:1211.5063}
  \item Adam: A Method for Stochastic Optimization (2014). \texttt{arXiv:1412.6980}
  \item Learning Phrase Representations using RNN Encoder-Decoder for
  Statistical Machine Translation (2014). \texttt{arXiv:1406.1078}
  \item Neural Machine Translation by Jointly Learning to Align and Translate (2014). \texttt{arXiv:1409.0473}
  \item Empirical Evaluation of Gated Recurrent Neural Networks on Sequence
  Modeling (2014). \texttt{arXiv:1412.3555}
  \item Deep Residual Learning for Image Recognition (2015). \texttt{arXiv:1512.03385}
  \item Unitary Evolution Recurrent Neural Networks (2015). \texttt{arXiv:1511.06464}
  \item Layer Normalization (2016). \texttt{arXiv:1607.06450}
  \item Gaussian Error Linear Units (GELUs) (2016). \texttt{arXiv:1606.08415}
  \item Pointer Sentinel Mixture Models (2016). \texttt{arXiv:1609.07843}
  \item Language Modeling with Gated Convolutional Networks (2016). \texttt{arXiv:1612.08083}
  \item The LAMBADA dataset: Word prediction requiring a broad discourse context (2016). \texttt{arXiv:1606.06031}
  \item Quasi-Recurrent Neural Networks (2016). \texttt{arXiv:1611.01576}
  \item Attention Is All You Need (2017). \texttt{arXiv:1706.03762}
  \item Decoupled Weight Decay Regularization (2017). \texttt{arXiv:1711.05101}
  \item Parallelizing Linear Recurrent Neural Nets Over Sequence Length (2017). \texttt{arXiv:1709.04057}
  \item BERT: Pre-training of Deep Bidirectional Transformers for Language
  Understanding (2018). \texttt{arXiv:1810.04805}
  \item GLUE: A Multi-Task Benchmark and Analysis Platform for Natural Language
  Understanding (2018). \texttt{arXiv:1804.07461}
  \item An Empirical Evaluation of Generic Convolutional and Recurrent Networks
  for Sequence Modeling (2018). \texttt{arXiv:1803.01271}
  \item Think you have Solved Question Answering? Try ARC, the AI2 Reasoning
  Challenge (2018). \texttt{arXiv:1803.05457}
  \item Speech Commands: A Dataset for Limited-Vocabulary Speech Recognition (2018). \texttt{arXiv:1804.03209}
  \item ListOps: A Diagnostic Dataset for Latent Tree Learning (2018). \texttt{arXiv:1804.06028}
  \item Can a Suit of Armor Conduct Electricity? A New Dataset for Open Book
  Question Answering (2018). \texttt{arXiv:1809.02789}
  \item Learning long-range spatial dependencies with horizontal gated-recurrent
  units (2018). \texttt{arXiv:1805.08315}
  \item Generating Long Sequences with Sparse Transformers (2019). \texttt{arXiv:1904.10509}
  \item Transformer-XL: Attentive Language Models Beyond a Fixed-Length Context (2019). \texttt{arXiv:1901.02860}
  \item RoBERTa: A Robustly Optimized BERT Pretraining Approach (2019). \texttt{arXiv:1907.11692}
  \item Exploring the Limits of Transfer Learning with a Unified Text-to-Text
  Transformer (2019). \texttt{arXiv:1910.10683}
  \item Compressive Transformers for Long-Range Sequence Modelling (2019). \texttt{arXiv:1911.05507}
  \item Megatron-LM: Training Multi-Billion Parameter Language Models Using
  Model Parallelism (2019). \texttt{arXiv:1909.08053}
  \item PyTorch: An Imperative Style, High-Performance Deep Learning Library (2019). \texttt{arXiv:1912.01703}
  \item PIQA: Reasoning about Physical Commonsense in Natural Language (2019). \texttt{arXiv:1911.11641}
  \item HellaSwag: Can a Machine Really Finish Your Sentence? (2019). \texttt{arXiv:1905.07830}
  \item WinoGrande: An Adversarial Winograd Schema Challenge at Scale (2019). \texttt{arXiv:1907.10641}
  \item Pay Less Attention with Lightweight and Dynamic Convolutions (2019). \texttt{arXiv:1901.10430}
  \item Transformer Dissection: A Unified Understanding of Transformer's
  Attention via the Lens of Kernel (2019). \texttt{arXiv:1908.11775}
  \item BoolQ: Exploring the Surprising Difficulty of Natural Yes/No Questions (2019). \texttt{arXiv:1905.10044}
  \item Transformers are RNNs: Fast Autoregressive Transformers with Linear
  Attention (2020). \texttt{arXiv:2006.16236}
  \item HiPPO: Recurrent Memory with Optimal Polynomial Projections (2020). \texttt{arXiv:2008.07669}
  \item Language Models are Few-Shot Learners (2020). \texttt{arXiv:2005.14165}
  \item Long Range Arena: A Benchmark for Efficient Transformers (2020). \texttt{arXiv:2011.04006}
  \item Reformer: The Efficient Transformer (2020). \texttt{arXiv:2001.04451}
  \item Longformer: The Long-Document Transformer (2020). \texttt{arXiv:2004.05150}
  \item Rethinking Attention with Performers (2020). \texttt{arXiv:2009.14794}
  \item Linformer: Self-Attention with Linear Complexity (2020). \texttt{arXiv:2006.04768}
  \item An Image is Worth 16x16 Words: Transformers for Image Recognition at
  Scale (2020). \texttt{arXiv:2010.11929}
  \item Big Bird: Transformers for Longer Sequences (2020). \texttt{arXiv:2007.14062}
  \item Efficient Content-Based Sparse Attention with Routing Transformers (2020). \texttt{arXiv:2003.05997}
  \item Efficient Transformers: A Survey (2020). \texttt{arXiv:2009.06732}
  \item Scaling Laws for Neural Language Models (2020). \texttt{arXiv:2001.08361}
  \item GLU Variants Improve Transformer (2020). \texttt{arXiv:2002.05202}
  \item Sparse Sinkhorn Attention (2020). \texttt{arXiv:2002.11296}
  \item Synthesizer: Rethinking Self-Attention in Transformer Models (2020). \texttt{arXiv:2005.00743}
  \item Efficiently Modeling Long Sequences with Structured State Spaces (2021). \texttt{arXiv:2111.00396}
  \item Combining Recurrent, Convolutional, and Continuous-time Models with
  Linear State-Space Layers (2021). \texttt{arXiv:2110.13985}
  \item CKConv: Continuous Kernel Convolution For Sequential Data (2021). \texttt{arXiv:2102.02611}
  \item RoFormer: Enhanced Transformer with Rotary Position Embedding (2021). \texttt{arXiv:2104.09864}
  \item The Pile: An 800GB Dataset of Diverse Text for Language Modeling (2021). \texttt{arXiv:2101.00027}
  \item Random Feature Attention (2021). \texttt{arXiv:2103.02143}
  \item Luna: Linear Unified Nested Attention (2021). \texttt{arXiv:2106.01540}
  \item On the Parameterization and Initialization of Diagonal State Space
  Models (2022). \texttt{arXiv:2206.11893}
  \item Diagonal State Spaces are as Effective as Structured State Spaces (2022). \texttt{arXiv:2203.14343}
  \item Simplified State Space Layers for Sequence Modeling (2022). \texttt{arXiv:2208.04933}
  \item Hungry Hungry Hippos: Towards Language Modeling with State Space Models (2022). \texttt{arXiv:2212.14052}
  \item FlashAttention: Fast and Memory-Efficient Exact Attention with
  IO-Awareness (2022). \texttt{arXiv:2205.14135}
  \item It's Raw! Audio Generation with State-Space Models (2022). \texttt{arXiv:2202.09729}
  \item How to Train Your HiPPO: State Space Models with Generalized Orthogonal
  Basis Projections (2022). \texttt{arXiv:2206.12037}
  \item Mega: Moving Average Equipped Gated Attention (2022). \texttt{arXiv:2209.10655}
  \item S4ND: Modeling Images and Videos as Multidimensional Signals Using State
  Spaces (2022). \texttt{arXiv:2210.06583}
  \item PaLM: Scaling Language Modeling with Pathways (2022). \texttt{arXiv:2204.02311}
  \item In-context Learning and Induction Heads (2022). \texttt{arXiv:2209.11895}
  \item Transformer Quality in Linear Time (2022). \texttt{arXiv:2202.10447}
  \item Long Range Language Modeling via Gated State Spaces (2022). \texttt{arXiv:2206.13947}
  \item What Makes Convolutional Models Great on Long Sequence Modeling? (2022). \texttt{arXiv:2210.09298}
  \item Convolutional State Space Models for Long-Range Spatiotemporal Modeling (2023). \texttt{arXiv:2310.19694}
  \item Spectral State Space Models (2023). \texttt{arXiv:2312.06837}
  \item Block-State Transformers (2023). \texttt{arXiv:2306.09539}
  \item Effectively Modeling Time Series with Simple Discrete State Spaces (2023). \texttt{arXiv:2303.09489}
  \item Simple Hardware-Efficient Long Convolutions for Sequence Modeling (2023). \texttt{arXiv:2302.06646}
  \item Mamba: Linear-Time Sequence Modeling with Selective State Spaces (2023). \texttt{arXiv:2312.00752}
  \item Structured State Space Models for In-Context Reinforcement Learning (2023). \texttt{arXiv:2303.03982}
  \item StableSSM: Alleviating the Curse of Memory in State-space Models through
  Stable Reparameterization (2023). \texttt{arXiv:2311.14495}
  \item Resurrecting Recurrent Neural Networks for Long Sequences (2023). \texttt{arXiv:2303.06349}
  \item State space models can express n-gram languages (2023). \texttt{arXiv:2306.17184}
  \item RWKV: Reinventing RNNs for the Transformer Era (2023). \texttt{arXiv:2305.13048}
  \item Hyena Hierarchy: Towards Larger Convolutional Language Models (2023). \texttt{arXiv:2302.10866}
  \item Retentive Network: A Successor to Transformer for Large Language Models (2023). \texttt{arXiv:2307.08621}
  \item Gated Linear Attention Transformers with Hardware-Efficient Training (2023). \texttt{arXiv:2312.06635}
  \item LLaMA: Open and Efficient Foundation Language Models (2023). \texttt{arXiv:2302.13971}
  \item GPT-4 Technical Report (2023). \texttt{arXiv:2303.08774}
  \item Llama 2: Open Foundation and Fine-Tuned Chat Models (2023). \texttt{arXiv:2307.09288}
  \item FlashAttention-2: Faster Attention with Better Parallelism and Work
  Partitioning (2023). \texttt{arXiv:2307.08691}
  \item Learning Long Sequences in Spiking Neural Networks (2024). \texttt{arXiv:2401.00955}
  \item State-Space Modeling in Long Sequence Processing: A Survey on Recurrence in the Transformer Era (2024). \texttt{arXiv:2406.09062}
  \item Mathematical Formalism for Memory Compression in Selective State Space
  Models (2024). \texttt{arXiv:2410.03158}
  \item Deep Learning-based Approaches for State Space Models: A Selective
  Review (2024). \texttt{arXiv:2412.11211}
  \item S7: Selective and Simplified State Space Layers for Sequence Modeling (2024). \texttt{arXiv:2410.03464}
  \item State Space Models as Foundation Models: A Control Theoretic Overview (2024). \texttt{arXiv:2403.16899}
  \item State Space Models on Temporal Graphs: A First-Principles Study (2024). \texttt{arXiv:2406.00943}
  \item Slot State Space Models (2024). \texttt{arXiv:2406.12272}
  \item Longhorn: State Space Models are Amortized Online Learners (2024). \texttt{arXiv:2407.14207}
  \item Oscillatory State-Space Models (2024). \texttt{arXiv:2410.03943}
  \item Transformers are SSMs: Generalized Models and Efficient Algorithms
  Through Structured State Space Duality (2024). \texttt{arXiv:2405.21060}
  \item Griffin: Mixing Gated Linear Recurrences with Local Attention for
  Efficient Language Models (2024). \texttt{arXiv:2402.19427}
  \item xLSTM: Extended Long Short-Term Memory (2024). \texttt{arXiv:2405.04517}
  \item Jamba: A Hybrid Transformer-Mamba Language Model (2024). \texttt{arXiv:2403.19887}
  \item Parallelizing Linear Transformers with the Delta Rule over Sequence
  Length (2024). \texttt{arXiv:2406.06484}
  \item The Illusion of State in State-Space Models (2024). \texttt{arXiv:2404.08819}
  \item MS-SSM: A Multi-Scale State Space Model for Efficient Sequence Modeling (2025). \texttt{arXiv:2512.23824}
  \item Advancing Intelligent Sequence Modeling: Evolution, Trade-offs, and Applications of State- Space Architectures from S4 to Mamba (2025). \texttt{arXiv:2503.18970}
  \item Geometric SSM: LTI State Space Models for Selective Tasks (2025). \texttt{arXiv:2505.17932}
  \item Message-Passing State-Space Models: Improving Graph Learning with Modern Sequence Modeling (2025). \texttt{arXiv:2505.18728}
  \item L2T-Hyena: Enhancing State-Space Models with an Adaptive Learn-to-Teach Framework (2025). \texttt{arXiv:2511.05926}
  \item Sequential-Parallel Duality in Prefix Scannable Models (2025). \texttt{arXiv:2506.10918}
  \item SeRpEnt: Selective Resampling for Expressive State Space Models (2025). \texttt{arXiv:2501.11729}
  \item Rethinking the long-range dependency in Mamba/SSM and transformer models (2025). \texttt{arXiv:2509.04226}
  \item Design Principles for Sequence Models via Coefficient Dynamics (2025). \texttt{arXiv:2510.09389}
  \item Numerical Investigation of Sequence Modeling Theory using Controllable Memory Functions (2025). \texttt{arXiv:2506.05678}
  \item Systolic Array-based Accelerator for Structured State-Space Models (2025). \texttt{arXiv:2507.21394}
  \item Time-Varying Deep State Space Models for Sequences with Switching Dynamics (2026). \texttt{arXiv:2605.15311}
\end{enumerate}

\end{document}
